\par\Needspace{12\baselineskip}\bigskip
\pdfbookmark[2]{Darstellung, Dritter Abschnitt}{pallas-II-section3}
\begin{center}
{\large\bfseries Dritter Abschnitt.\par}
\medskip
\textit{Methode zur Berechnung der endlichen Änderungen der Elemente während eines begrenzten Zeitraums.}
\end{center}
\pdfbookmark[3]{Darstellung, 17}{pallas-II-article17}
\begin{center}17.\end{center}
\noindent\hspace*{\parindent}Die Kleinheit der störenden Massen im Vergleich zur Masse der Sonne ist ein äußerst günstiger Umstand in der Theorie der Planetenstörungen. Genau genommen müsste man bei der Berechnung der \SourceLeaf{467}momentanen Änderungen jedes Elements die wahren Werte der Elemente verwenden, die in den Ausdruck für diese Änderungen eingehen. Angesichts der Kleinheit dieser Änderungen selbst wird das Ergebnis jedoch nur sehr wenig abweichen, wenn man lediglich Näherungswerte der Elemente verwendet. Zumindest erhält man auf diese Weise eine erste Näherung, die sich später verbessern lässt, wenn man es für nötig hält.
\par Zunächst bestimmt man also nach den bekannten Methoden rein elliptische Elemente, die mit den Beobachtungen möglichst gut übereinstimmen. Geht es nun nur um die Theorie der Bewegung des Planeten während eines begrenzten Zeitraums, so teilt man diesen zunächst in mehrere gleiche Intervalle, für die wir jeweils \(\theta\) Tage ansetzen. Zweckmäßig wählt man die erste Epoche dieser Intervalle, die wir mit \(T\) bezeichnen, etwas vor dem Zeitraum, für den man die Bewegung berechnen will, und setzt die Reihe \(T,\ T+\theta,\ T+2\theta,\ T+3\theta\) usw. noch etwas über diesen Zeitraum hinaus fort. Nach den oben entwickelten Formeln und anhand der Näherungswerte der Elemente berechnet man die Zahlenwerte der momentanen Änderungen der verschiedenen Elemente zu den Epochen \(T,\ T+\theta,\ T+2\theta\) usw., nämlich die Werte von \(\dfrac{\mathrm{d}i}{\mathrm{d}t},\ \dfrac{\mathrm{d}\AscendingNode}{\mathrm{d}t},\ \dfrac{\mathrm{d}n}{\mathrm{d}t},\ \dfrac{\mathrm{d}\varphi}{\mathrm{d}t},\ \dfrac{\mathrm{d}\tilde\omega}{\mathrm{d}t},\ \dfrac{\mathrm{d}\epsilon}{\mathrm{d}t}\). Den Wert von \(\dfrac{\mathrm{d}a}{\mathrm{d}t}\) benötigt man nicht, da der Wert von \(a\) stets mit dem von \(n\) verknüpft ist. Es bleibt also, diese sechs Differentiale zu integrieren und durch eine zweite Integration die mittlere Bewegung \(\int n\mathrm{d}t\) zu bestimmen. Die folgende Methode erscheint uns hierfür am bequemsten.
